\begin{figure}[htbp]
	\centering
	\begin{tikzpicture}[font=\small,
		lembo/.style={trvconcio, fill=black!5}]
		\def\Lc{1.15}\def\Hc{0.72}\def\g{0.5}
		% ascisse comuni ai tre pannelli: C^- in [0, Lc], C^+ in [Lc+g, 2Lc+g]
		\pgfmathsetmacro{\xR}{\Lc}
		\pgfmathsetmacro{\xL}{\Lc+\g}
		\pgfmathsetmacro{\xE}{2*\Lc+\g}
		\pgfmathsetmacro{\xM}{\Lc+\g/2}
		% ===============================================================
		% (a) distorsione estensionale e forza normale
		% ===============================================================
		\begin{scope}
			\draw[lembo] (0,-\Hc/2) rectangle (\xR,\Hc/2);
			\draw[lembo] (\xL,-\Hc/2) rectangle (\xE,\Hc/2);
			\node[above] at (\Lc/2,\Hc/2) {$\mathcal{C}^-$};
			\node[above] at (\xL+\Lc/2,\Hc/2) {$\mathcal{C}^+$};
			\draw[trvforza] (\xR-0.5,0) -- (\xR,0);
			\node[above=2pt] at (\xR-0.3,0.02) {$N$};
			\draw[trvforza] (\xL+0.5,0) -- (\xL,0);
			\node[above=2pt] at (\xL+0.3,0.02) {$N$};
			\draw[trvquota, black!55] (\xR,-\Hc/2-0.25) -- node[below=1pt, black!70] {$\varrho_w$} (\xL,-\Hc/2-0.25);
			\node at (\xM,-1.35) {(a) estensionale};
		\end{scope}
		% ===============================================================
		% (b) distorsione trasversale e taglio
		% ===============================================================
		\begin{scope}[xshift=3.75cm]
			\def\dv{0.42}
			\draw[lembo] (0,-\Hc/2) rectangle (\xR,\Hc/2);
			\draw[lembo] (\xL,-\Hc/2-\dv) rectangle (\xE,\Hc/2-\dv);
			\node[above] at (\Lc/2,\Hc/2) {$\mathcal{C}^-$};
			\node[below] at (\xL+\Lc/2,-\Hc/2-\dv) {$\mathcal{C}^+$};
			\draw[trvcostruzione] (\xR,\Hc/2) -- (\xL,\Hc/2);
			\draw[trvquota, black!55] (\xM,\Hc/2) -- node[right=1pt, black!70] {$\varrho_v$} (\xM,\Hc/2-\dv);
			\draw[trvforza] (\xR-0.12,0.22) -- (\xR-0.12,-0.22);
			\node[left] at (\xR-0.12,0) {$T$};
			\draw[trvforza] (\xL+0.12,-\dv-0.22) -- (\xL+0.12,-\dv+0.22);
			\node[right] at (\xL+0.12,-\dv) {$T$};
			\node at (\xM,-1.35) {(b) trasversale};
		\end{scope}
		% ===============================================================
		% (c) distorsione rotazionale e momento flettente
		% ===============================================================
		\begin{scope}[xshift=7.5cm]
			\def\t{15}
			\draw[lembo] (0,-\Hc/2) rectangle (\xR,\Hc/2);
			\begin{scope}[rotate around={\t:(\xL,0)}]
				\draw[lembo] (\xL,-\Hc/2) rectangle (\xE,\Hc/2);
				\draw[line width=.5pt, black!45] (\xL-0.35,0) -- (\xE+0.45,0);
				\draw[trvcoppia] (\xL+0.06,0.2) arc[start angle=90, end angle=270, radius=0.2];
			\end{scope}
			\draw[line width=.5pt, black!45] (-0.1,0) -- (\xE+0.45,0);
			\node[above] at (\Lc/2,\Hc/2) {$\mathcal{C}^-$};
			\node[above] at (\xL+\Lc/2-0.1,\Hc/2+0.2) {$\mathcal{C}^+$};
			\draw[trvcoppia] (\xR-0.06,-0.2) arc[start angle=-90, end angle=90, radius=0.2];
			\node[below] at (\xR,-\Hc/2) {$M$};
			\node[below] at ($(\xL,0)+(\t-90:\Hc/2)$) {$M$};
			\draw[trvdist] (\xL+1.45,0) arc[start angle=0, end angle=\t, radius=1.45];
			\node[right] at (\xL+1.47,0.18) {$\varrho_\vartheta$};
			\node at (\xM,-1.35) {(c) rotazionale};
		\end{scope}
	\end{tikzpicture}
	\caption{Distorsioni e caratteristiche della sollecitazione
		duali: (a)~$\varrho_w$ e $N$; (b)~$\varrho_v$ e $T$;
		(c)~$\varrho_\vartheta$ e $M$. Le caratteristiche
		agiscono con versi opposti sulle due facce: la coppia
		$(\bm{\sigma}, \bm{\varrho})$ entra nell'identità
		bilineare con segno discorde,
		come nella~\eqref{eq:lavoro_continuita}.}
	\label{fig:distorsioni_concio}
\end{figure}
